\documentclass[11pt,a4paper]{article}
\usepackage[margin=27mm]{geometry}
\usepackage{amsmath,amssymb,amsthm,hyperref}
\newtheorem{theorem}{Theorem}
\newcommand{\gen}[1]{\langle#1\rangle}
\newcommand{\Sz}{\operatorname{Sz}}
\title{A second failure of the thinness assertion}
\author{Human Compression project}\date{}
\begin{document}\maketitle
\begin{abstract}
The thinness clause for prime $q-1>7$ is separate from the clause for
$q=8$. We settle it by an explicit example in $\Sz(32)$, using the same
coset argument as for the previously verified $\Sz(8)$ example. Its
parabolics have orders $62,62,82$, and each actual panel has three
elements although the corresponding algebraic residue has two.
\end{abstract}

Work in $F=\mathbb F_2[X]/(X^5+X^2+1)$. Integers from zero to31 encode
coefficients in the basis $1,X,X^2,X^3,X^4$; field addition is
exclusive-or. The independent arithmetic reader checks all31 nonzero
polynomial classes have inverses. In the classical Suzuki model put
\[
 a=\begin{pmatrix}0&0&0&1\\0&0&1&0\\0&1&0&0\\1&0&0&0\end{pmatrix},\quad
 b=\begin{pmatrix}12&3&0&1\\3&0&1&0\\20&4&0&3\\30&20&3&12\end{pmatrix},\quad
 c=\begin{pmatrix}26&5&0&1\\5&0&1&0\\29&16&0&5\\19&29&5&26\end{pmatrix}.
\]
Their membership in the ambient group is explicit. With $w=a$ and
$z(t)=u(0,t)$, where
\[
 u(s,t)=\begin{pmatrix}
 1&0&0&0\\s&1&0&0\\s^9+t&s^8&1&0\\s^{10}+st+t^8&t&s&1
 \end{pmatrix},
\]
we have $b=z(3)wz(3)$ and
$c=u(1,5)wz(1)wu(1,5)^{-1}$. The ambient matrix construction and
maximality of the full split normalizer $D_{62}$ in $\Sz(32)$ remain
explicit classical premises.

Literal polynomial multiplication verifies that $a,b,c$ are distinct
nonidentity involutions and that
\[
 |ab|=31,\qquad |ac|=31,\qquad |bc|=41,\qquad
 (ac)^{15}(ab)^{15}a=(bc)^{20}.                         \tag{1}
\]
Define $G_0=\gen{a,b}$, $G_1=\gen{a,c}$ and $G_2=\gen{b,c}$.
Their respective orders are $62,62,82$. Coprimality gives
$G_0\cap G_2=\gen b$ and $G_1\cap G_2=\gen c$.
The subgroup $G_0\cap G_1$ contains $\gen a$, which is maximal in
the prime dihedral group $G_0$. If the intersection were larger,
then $G_0=G_1$, placing $b,c$ in a group of order62 and contradicting
$|bc|=41$. Thus $G_0\cap G_1=\gen a$, and the triple intersection is
trivial.

\begin{theorem}
The rank-three coset geometry $\Gamma(\Sz(32);G_0,G_1,G_2)$ is RWPRI
and satisfies the literal algebraic local $2$-transitivity condition
of L114. It is not thin and is not flag transitive.
\end{theorem}
\begin{proof}
An incident pair of left cosets has a common representative and hence
extends to a chamber. Thus the pregeometry is a geometry.
For the empty subset in the RWPRI definition, $G_0$ is maximal in
$G=\Sz(32)$. Each pair intersection has prime index31 or41 in the
corresponding parabolic, and the trivial triple intersection is maximal
in each pair intersection. These verify every required subset.
Each source algebraic rank-one residue is the regular two-point
action $C_2/1$, so is $2$-transitive.

The base edge $(G_1,G_2)$ has the standard extensions $G_0,cG_0$.
Put $x=(ac)^{15}$. The coset $xG_0$ meets $G_1$ at $x$ and, by (1),
meets $G_2$ at $(bc)^{20}$. It differs from $G_0$, since $x$ has
order31 and $G_0\cap G_1=\gen a$. If $xG_0=cG_0$, then
$cx\in\gen a$. The alternative $cx=1$ contradicts the orders of
$c,x$. The alternative $cx=a$ gives $(ac)^{15}=ca=(ac)^{30}$,
contradicting $|ac|=31$. We have found a third chamber, so the geometry
is not thin. This chamber has no common representative: such a
representative would lie in $G_1\cap G_2=\{1,c\}$ and identify its
third coset with one of the standard cosets. It therefore cannot be a
translate of the base chamber, proving failure of flag transitivity.
\end{proof}

The independent reader checks every product and inverse in the three
returned subgroup lists and verifies generation, intersections and all
cosets in each panel. It obtains exactly three cosets in each case,
with the following representatives:
\[
\begin{array}{c|c|c}
\text{missing type}&\text{standard representatives}&\text{extra representative}\\\hline
0&1,c&(ac)^{15}\\
1&1,b&(ab)^{15}\\
2&1,a&(ba)^{15}
\end{array}
\]
The reader explicitly matches each displayed representative with its
complete returned coset. Thus the table is certified along with (1).

\paragraph{Source scope and attribution.}
The source is D.~Leemans, \emph{The rank 3 geometries of the simple
Suzuki groups $\Sz(q)$}, Note di Matematica 19(1) (1999), 43--63.
Its definitions on printed pp.~45--46 use algebraic residues formed
from parabolic intersections, and count actual chambers for thinness.
The example has $q-1=31>7$ prime and therefore refutes the thinness
clause in the second branch of Theorem~6.1 on printed p.~61.
The previously accepted rank bound and rank-two count in that branch
are unaffected. A version requiring flag transitivity or local
$2$-transitivity on full actual panels has different hypotheses.

The field and split pair were retained from earlier work. The new
contribution is this third involution, its short coset identity and
the exact source application, reusing the preceding $\Sz(8)$ argument.
The bounded discovery stopped after156 prescribed candidates. Its
unaltered output is retained; the independent reader verifies this
witness alone and makes no classification claim. No new original,
principal-refutation count, open-conjecture, Lean or compression credit
is assigned. Distinct mathematical review is required before adoption.
\end{document}
